\chapter{Introduction}
\label{chap:intro}

Modern image classifiers are trained on data that no single party can inspect,
which turns the training set into an attack surface. In a backdoor poisoning
attack, an adversary adds a trigger, such as a small patch or a faint blended
pattern, to a few percent of the training images and relabels
them~\cite{gu2019badnets,chen2017blended}. The resulting model behaves normally
on clean inputs but predicts the target class whenever the trigger is
present~\cite{li2024backdoorsurvey}.

Institutions that train a joint model on confidential data can use secure
multi-party computation (MPC), in which the data stays secret-shared and no
party sees it~\cite{mohassel2018aby3,knott2021crypten}. MPC protects the
confidentiality of the inputs but not their honesty, so a poisoning defense has
to run inside MPC itself. The strongest dataset
purification defenses, however, first train a network on the poisoned data and
then apply singular value decomposition or clustering to its
features~\cite{tran2018spectral,chen2019activation}, and both steps are
prohibitively expensive under MPC.

Perceptual hash functions offer a cheap alternative. A perceptual hash maps an
image to a short bit string so that similar images obtain similar
strings~\cite{zauner2010phash}, and most classical constructions consist of a
public linear transform, free under secret sharing, followed by a
binarisation. Yet perceptual hashes have been studied in
machine learning security only as the target of
attacks~\cite{prokos2023squint}, and it is unknown whether they can detect
poisoned samples, which triggers they miss and what they cost inside MPC. This thesis therefore addresses three research questions:
\begin{enumerate}[label=\textbf{RQ\arabic*}, leftmargin=*, nosep]
  \item How effective are perceptual hashes, combined with statistics computed
    on their bits, at removing samples poisoned by state-of-the-art backdoor
    attacks?
  \item Which properties of a trigger determine whether it can be detected from
    per-sample hash bits?
  \item What does hash-based purification cost under MPC protocols compared with
    model-based defenses?
\end{enumerate}

To answer them, the thesis writes classical perceptual hashes in a common
MPC-friendly form, proposes a family of residual hashes that keep the
high-frequency component in which invisible triggers reside, and evaluates
26 hashes and seven detection statistics against 13 attacks on
CIFAR-10~\cite{krizhevsky2009cifar}, measuring their cost under three-party and
two-party MPC protocols. The study found that classical hashes detected only visible patches,
whereas residual hashes detected every sample-agnostic trigger tested and
reduced the attack success rate of a ResNet-18 from 1.00 to 0.00 for an adaptive
attack and to 0.15 for a blended one; partial removal, however, left the BadNets
backdoor intact, and sample-specific triggers remained undetectable. Purifying
the 50\,000 CIFAR-10 training images cost about 29~GB of communication and 78~s
on a local network, about $2\cdot10^{9}$ times less than a model-based defense.

Chapter~\ref{chap:lr} reviews related work, Chapter~\ref{chap:met} describes
the methodology, Chapter~\ref{chap:impl} the implementation and
Chapter~\ref{chap:eval} the results; Chapter~\ref{chap:conclusion} concludes.
